% preamble.tex -- shared preamble for the combined monograph.
% Chapter-specific macros live inside each chapter, wrapped in \begin{chaptermacros}...\end{chaptermacros}
% or defined with \chapmacro inside the chapter's local group (see CONVENTIONS.md).

\usepackage[T1]{fontenc}
\usepackage[utf8]{inputenc}
\IfFileExists{lmodern.sty}{\usepackage{lmodern}}{}
\usepackage[a4paper,margin=1.15in]{geometry}
\usepackage{amsmath,amssymb,amsfonts,mathtools,amsthm}
\usepackage{mathrsfs,bm,cancel,xfrac}
\usepackage{tikz-cd}
\usepackage{xcolor}
\usepackage{graphicx}
\usepackage{booktabs,array,longtable,multirow}
\usepackage[shortlabels]{enumitem}
\usepackage[numbers,sort&compress]{natbib}
\usepackage[hypertexnames=false,bookmarksnumbered=true]{hyperref}
\hypersetup{colorlinks=true,linkcolor=blue!60!black,citecolor=green!40!black,urlcolor=blue!60!black,
  pdftitle={Continuity, Testing Regions and Fragility of Quantum Channel Divergences},pdfauthor={M. Jones}}
\usepackage{bookmark}

\graphicspath{{figures/}}
\emergencystretch=2em
\allowdisplaybreaks[2]

% ---------- numbering: chapter.section.number ----------
\numberwithin{equation}{section}

\theoremstyle{plain}
\newtheorem{theorem}{Theorem}[section]
\newtheorem{lemma}[theorem]{Lemma}
\newtheorem{proposition}[theorem]{Proposition}
\newtheorem{corollary}[theorem]{Corollary}
\newtheorem{conjecture}[theorem]{Conjecture}
\newtheorem{claim}[theorem]{Claim}
\newtheorem{fact}[theorem]{Fact}
\newtheorem{assumption}[theorem]{Assumption}
\newtheorem*{theorem*}{Theorem}
\newtheorem*{lemma*}{Lemma}
\newtheorem*{proposition*}{Proposition}
\newtheorem*{corollary*}{Corollary}
\newtheorem*{conjecture*}{Conjecture}
\newtheorem*{claim*}{Claim}

\theoremstyle{definition}
\newtheorem{definition}[theorem]{Definition}
\newtheorem{example}[theorem]{Example}
\newtheorem{problem}[theorem]{Problem}
\newtheorem{question}[theorem]{Question}
\newtheorem*{definition*}{Definition}
\newtheorem*{example*}{Example}
\newtheorem*{question*}{Question}

\theoremstyle{remark}
\newtheorem{remark}[theorem]{Remark}
\newtheorem{notation}[theorem]{Notation}
\newtheorem{observation}[theorem]{Observation}
\newtheorem*{remark*}{Remark}
\newtheorem*{notation*}{Notation}

% ---------- global notation (used in every chapter; see chapters/notation.tex) ----------
\newcommand{\tr}{\operatorname{tr}}
\newcommand{\Tr}{\operatorname{Tr}}
\newcommand{\id}{\mathrm{id}}
\newcommand{\vectorise}{\operatorname{vec}}
\DeclareMathOperator{\range}{range}
\providecommand{\rank}{\operatorname{rank}}

% ---------- chapter-local macro helper ----------
% \chapmacro{\foo}{definition} or \chapmacro[1]{\foo}{definition with #1}
% (re)defines \foo whether or not it already exists. Use only inside a chapter's local group.
\newcommand{\chapmacro}[3][0]{\providecommand{#2}{}\renewcommand{#2}[#1]{#3}}
\newcommand{\chapmathop}[2]{\providecommand{#1}{}\renewcommand{#1}{\operatorname{#2}}}

% ---------- chapter summary block ----------
\newenvironment{chaptersummary}{\begin{quote}\small\noindent\textbf{Summary.}\ }{\end{quote}\medskip}
